% ============================================================
% Chapter 5: Human Factors, Social Engineering and Digital Safety
% Language: Greek — edit freely; regenerate from src/content if preferred.
% ============================================================
\chapter{Ανθρώπινος Παράγοντας, Κοινωνική Μηχανική και Ψηφιακή Ασφάλεια}\label{ch:5}
\chaptermeta{Ψυχολογία της πειθούς · Ηλεκτρονικό ψάρεμα και οι παραλλαγές του · Φυσικά κανάλια · Stalkerware και μεταδεδομένα · Παραβίαση επιχειρηματικού email · Κουλτούρα ασφάλειας}{Επίπεδο: Εισαγωγικό επίπεδο · Ανθρωποκεντρική προσέγγιση}{Χρόνος μελέτης: 5–7 ώρες}

Τα τεχνικά μέτρα προστατεύουν τα συστήματα, όμως τα συστήματα τα χειρίζονται άνθρωποι, και οι άνθρωποι μπορούν να πειστούν. Η κοινωνική μηχανική εκμεταλλεύεται την εμπιστοσύνη, την προθυμία για βοήθεια, τον φόβο και τη συνήθεια, ώστε ένας νόμιμος χρήστης να εκτελέσει μια ενέργεια που ωφελεί τον επιτιθέμενο —να ανοίξει ένα συνημμένο, να αποκαλύψει έναν κωδικό, να εγκρίνει μια πληρωμή. Οι κλαδικές αναφορές παραβιάσεων διαπιστώνουν σταθερά ότι ο ανθρώπινος παράγοντας εμπλέκεται στην πλειονότητα των επιτυχημένων εισβολών.

Το κεφάλαιο εξηγεί γιατί λειτουργεί η κοινωνική μηχανική, παρουσιάζει τις κύριες μεθόδους επίθεσης και εξετάζει δύο απειλές που έχουν αυξηθεί σημαντικά τα τελευταία χρόνια: την παραβίαση επιχειρηματικού ηλεκτρονικού ταχυδρομείου και την «κόπωση» του ελέγχου ταυτότητας πολλαπλών παραγόντων. Εξετάζει επίσης την προσωπική ψηφιακή ασφάλεια, συμπεριλαμβανομένου του λογισμικού παρακολούθησης (stalkerware) και της διαρροής μεταδεδομένων. Το κεφάλαιο κλείνει με ένα εποικοδομητικό μήνυμα: οι χρήστες δεν είναι ο «πιο αδύναμος κρίκος» που πρέπει να κατηγορείται, αλλά ένα επίπεδο ανίχνευσης που μπορεί να ενισχυθεί μέσω καλού σχεδιασμού και μιας θετικής κουλτούρας αναφοράς.

\begin{outcomesbox}{Μαθησιακά αποτελέσματα}
Μετά τη μελέτη του κεφαλαίου θα πρέπει να μπορείτε:
\begin{itemize}
  \item Να εξηγείτε τις ψυχολογικές αρχές που εκμεταλλεύονται οι κοινωνικοί μηχανικοί.
  \item Να διακρίνετε το ηλεκτρονικό ψάρεμα, το στοχευμένο ψάρεμα, το whaling, το vishing, το smishing, την κατασκευή προσχήματος και το δόλωμα.
  \item Να περιγράφετε τα φυσικά κανάλια επίθεσης και τα κανάλια παρατήρησης, καθώς και τα αντίμετρά τους.
  \item Να αναλύετε μια παραβίαση επιχειρηματικού email και μια επίθεση κόπωσης MFA και να προτείνετε πολυεπίπεδες άμυνες.
  \item Να σχεδιάζετε μέτρα ευαισθητοποίησης και δείκτες που μειώνουν την ευαλωτότητα χωρίς να επιρρίπτουν ευθύνες στους χρήστες.
\end{itemize}
\end{outcomesbox}

\section{Η ψυχολογία της εκμετάλλευσης}\label{sec:5.1}
Η κοινωνική μηχανική επιτυγχάνει επειδή στοχεύει νοητικές συντομεύσεις που υπό κανονικές συνθήκες μάς εξυπηρετούν. Ο ψυχολόγος Robert Cialdini περιέγραψε ορισμένες \textbf{αρχές της πειθούς} που οι επιτιθέμενοι εκμεταλλεύονται συστηματικά. \emph{Αυθεντία}: οι άνθρωποι τείνουν να υπακούουν σε οδηγίες που φαίνεται να προέρχονται από έναν προϊστάμενο, το τμήμα πληροφορικής ή την αστυνομία. \emph{Επείγον και σπανιότητα}: μια προθεσμία («ο λογαριασμός σας θα κλείσει σε μία ώρα») μειώνει τον διαθέσιμο χρόνο για προσεκτική σκέψη. \emph{Κοινωνική απόδειξη}: υποθέτουμε ότι μια ενέργεια είναι ασφαλής αν φαίνεται ότι την κάνουν και άλλοι. \emph{Συμπάθεια} και \emph{αμοιβαιότητα}: είμαστε πιο πρόθυμοι να βοηθήσουμε κάποιον που είναι φιλικός ή που μας έχει κάνει μια μικρή χάρη.

Οι αρχές αυτές είναι ισχυρές επειδή παρακάμπτουν τη συνειδητή σκέψη και ενεργοποιούν γρήγορες, αυτόματες αντιδράσεις. Το άγχος, η κούραση και ο υπερβολικός όγκος πληροφοριών καθιστούν όλους πιο ευάλωτους, ανεξάρτητα από την ευφυΐα ή τις τεχνικές τους γνώσεις. Οι αποτελεσματικές άμυνες δεν βασίζονται, επομένως, στη διαρκή εγρήγορση των ανθρώπων· εισάγουν \emph{τριβή την κατάλληλη στιγμή} —για παράδειγμα, υποχρεωτική επιβεβαιωτική κλήση πριν από οποιαδήποτε αλλαγή τραπεζικών στοιχείων— ώστε η αυτόματη αντίδραση να διακόπτεται ακριβώς όταν έχει σημασία.

\section{Μέθοδοι επίθεσης: από το μαζικό ψάρεμα έως την κατασκευή προσχήματος}\label{sec:5.2}
Το \textbf{ηλεκτρονικό ψάρεμα} (phishing) είναι η μαζική αποστολή δόλιων μηνυμάτων που μιμούνται αξιόπιστους οργανισμούς, με σκοπό την υποκλοπή διαπιστευτηρίων ή την εγκατάσταση κακόβουλου λογισμικού. Το \textbf{στοχευμένο ψάρεμα} (spear-phishing) απευθύνεται σε συγκεκριμένο πρόσωπο ή ομάδα και χρησιμοποιεί στοιχεία που έχουν συλλεχθεί με έρευνα —ονόματα συναδέλφων, τρέχοντα έργα— ώστε να φαίνεται αξιόπιστο. Το \textbf{whaling} είναι στοχευμένο ψάρεμα που απευθύνεται σε ανώτατα στελέχη. Οι ίδιες τεχνικές εφαρμόζονται και σε άλλα κανάλια: το \textbf{vishing} χρησιμοποιεί τηλεφωνικές κλήσεις, συχνά με πλαστογραφημένη αναγνώριση καλούντος και, όλο και περισσότερο, με φωνές που παράγονται από τεχνητή νοημοσύνη, ενώ το \textbf{smishing} χρησιμοποιεί SMS ή εφαρμογές ανταλλαγής μηνυμάτων.

Η \textbf{κατασκευή προσχήματος} (pretexting) συνίσταται στην επινόηση ενός αληθοφανούς σεναρίου —ένας ελεγκτής που χρειάζεται μια αναφορά, ένας τεχνικός που πρέπει να «επαληθεύσει» έναν κωδικό— για την απόκτηση πληροφοριών ή πρόσβασης. Το \textbf{δόλωμα} (baiting) αφήνει ένα ελκυστικό αντικείμενο, όπως ένα USB με την ετικέτα «Μισθοδοσία 2026», σε σημείο όπου το θύμα θα το βρει και θα το συνδέσει στον υπολογιστή του. Το κοινό στοιχείο είναι ότι ο επιτιθέμενος δημιουργεί μια κατάσταση στην οποία η επιβλαβής ενέργεια μοιάζει φυσιολογική, εξυπηρετική ή επείγουσα.

\section{Φυσικά κανάλια και κανάλια παρατήρησης}\label{sec:5.3}
Δεν πραγματοποιούνται όλες οι επιθέσεις κοινωνικής μηχανικής στο διαδίκτυο. Το \textbf{tailgating} (ή \emph{piggybacking}) σημαίνει να ακολουθεί κανείς ένα εξουσιοδοτημένο πρόσωπο μέσα από μια ασφαλισμένη πόρτα, συχνά κρατώντας κούτες, ώστε το να του κρατήσει κάποιος την πόρτα να φαίνεται απλή ευγένεια. Η \textbf{παρακολούθηση πάνω από τον ώμο} (shoulder surfing) είναι η παρατήρηση μιας οθόνης ή ενός πληκτρολογίου για την υποκλοπή κωδικών ή PIN, ενώ η \textbf{έρευνα στα απορρίμματα} (dumpster diving) είναι η αναζήτηση χρήσιμων πληροφοριών σε πεταμένα έγγραφα και εξοπλισμό. Η φυσική πρόσβαση είναι ιδιαίτερα επικίνδυνη, επειδή μπορεί να παρακάμψει εντελώς πολλά δικτυακά μέτρα —ένας επιτιθέμενος που φτάνει σε έναν ξεκλείδωτο σταθμό εργασίας ή σε μια δικτυακή θύρα βρίσκεται ήδη «μέσα».

Τα αντίμετρα συνδυάζουν τεχνολογία και συμπεριφορά: θαλάμους διπλής πόρτας (mantraps) και περιστροφικές πύλες που επιτρέπουν τη διέλευση ενός ατόμου κάθε φορά, εμφανείς κάρτες επισκεπτών, πολιτική καθαρού γραφείου και κλειδωμένης οθόνης, φίλτρα απορρήτου στις οθόνες, καταστροφή εγγράφων σε τεμαχιστές εγκάρσιας κοπής και πιστοποιημένη καταστροφή αποθηκευτικών μέσων. Το προσωπικό πρέπει να αισθάνεται ότι δικαιούται να ζητήσει ευγενικά στοιχεία από ένα άγνωστο πρόσωπο· αυτό γίνεται ευκολότερο όταν η διοίκηση το υποστηρίζει ρητά.

\section{Προσωπική ψηφιακή ασφάλεια: stalkerware και διαρροή μεταδεδομένων}\label{sec:5.4}
Ορισμένες από τις πιο επιβλαβείς επιθέσεις δεν πραγματοποιούνται από απομακρυσμένους εγκληματίες, αλλά από πρόσωπα του στενού περιβάλλοντος του θύματος. Το \textbf{stalkerware} είναι εμπορικά διαθέσιμο λογισμικό το οποίο, μόλις εγκατασταθεί σε ένα κινητό τηλέφωνο, προωθεί κρυφά μηνύματα, τοποθεσία, φωτογραφίες και ιστορικό κλήσεων σε άλλο πρόσωπο. Συνδέεται συχνά με την κακοποίηση από σύντροφο. Προειδοποιητικά σημάδια αποτελούν η ανεξήγητη εξάντληση της μπαταρίας, άγνωστες εφαρμογές με δικαιώματα διαχειριστή συσκευής και ένας σύντροφος που γνωρίζει πράγματα που δεν θα έπρεπε. Επειδή η αφαίρεση του stalkerware μπορεί να ειδοποιήσει τον θύτη και να αυξήσει τον κίνδυνο, τα θύματα πρέπει να αναζητούν υποστήριξη από εξειδικευμένους φορείς πριν προβούν σε οποιαδήποτε ενέργεια.

Τα \textbf{μεταδεδομένα} —δεδομένα για τα δεδομένα— μπορεί να αποκαλύπτουν πολύ περισσότερα από όσα σκοπεύαμε. Τα μεταδεδομένα EXIF μιας φωτογραφίας μπορεί να περιέχουν τις ακριβείς συντεταγμένες GPS του σημείου λήψης, το μοντέλο της συσκευής και την ώρα· τα έγγραφα γραφείου μπορεί να περιλαμβάνουν το όνομα του συντάκτη και το ιστορικό αναθεωρήσεων. Πριν από τη δημόσια κοινοποίηση αρχείων, οι χρήστες πρέπει να αφαιρούν αυτά τα μεταδεδομένα, για παράδειγμα με την εντολή \texttt{exiftool -all=} ή με τις ενσωματωμένες λειτουργίες «κατάργησης ιδιοτήτων» των λειτουργικών συστημάτων και των σουιτών γραφείου.

\section{Παραβίαση επιχειρηματικού email και κόπωση MFA}\label{sec:5.5}
Η \textbf{παραβίαση επιχειρηματικού ηλεκτρονικού ταχυδρομείου} (Business Email Compromise, BEC) είναι μια απάτη κατά την οποία ο επιτιθέμενος υποδύεται ένα ανώτατο στέλεχος, έναν προμηθευτή ή έναν δικηγόρο —είτε παραβιάζοντας ένα πραγματικό γραμματοκιβώτιο είτε καταχωρίζοντας ένα παρόμοιο όνομα τομέα— και ζητά μια επείγουσα πληρωμή ή αλλαγή τραπεζικών στοιχείων. Η BEC συχνά δεν περιλαμβάνει καθόλου κακόβουλο λογισμικό και, επομένως, διαφεύγει από τα τεχνικά φίλτρα· βασίζεται αποκλειστικά στην αυθεντία και στο επείγον. Οι αποτελεσματικές άμυνες είναι τόσο διαδικαστικές όσο και τεχνικές: οι αλλαγές στοιχείων πληρωμής επαληθεύονται μέσω γνωστού αριθμού τηλεφώνου, οι μεταφορές υψηλής αξίας απαιτούν έγκριση από δύο πρόσωπα (\emph{δημιουργός–ελεγκτής}) και ο έλεγχος αυθεντικότητας email (SPF, DKIM και πολιτική DMARC \texttt{p=reject}) δυσχεραίνει σημαντικά την πλαστογράφηση τομέα.

Η \textbf{κόπωση MFA} (MFA fatigue ή \emph{push bombing}) στοχεύει χρήστες που επαληθεύουν την ταυτότητά τους μέσω ειδοποιήσεων push. Έχοντας υποκλέψει έναν κωδικό, ο επιτιθέμενος προκαλεί δεκάδες αιτήματα έγκρισης σύνδεσης, συχνά αργά τη νύχτα, έως ότου ο εξαντλημένος χρήστης πατήσει «Έγκριση» —μερικές φορές μετά από τηλεφώνημα ενός ψεύτικου γραφείου υποστήριξης. Τα αντίμετρα περιλαμβάνουν την \emph{αντιστοίχιση αριθμού} (ο χρήστης πρέπει να πληκτρολογήσει έναν κωδικό που εμφανίζεται στην οθόνη σύνδεσης), όρια στον αριθμό των αιτημάτων και μεθόδους ανθεκτικές στο ψάρεμα, όπως τα κλειδιά ασφαλείας FIDO2 (Κεφάλαιο 8).

\section{Οικοδόμηση κουλτούρας ασφάλειας που ενθαρρύνει την αναφορά}\label{sec:5.6}
Τα προγράμματα ευαισθητοποίησης συχνά μετρούν το λάθος πράγμα. Το \textbf{ποσοστό κλικ} στις προσομοιωμένες εκστρατείες ηλεκτρονικού ψαρέματος αντανακλά κυρίως τη δυσκολία της προσομοίωσης· μπορεί να αυξηθεί ή να μειωθεί κατά βούληση. Πολύ πιο ουσιαστικός δείκτης είναι το \textbf{ποσοστό αναφοράς} —το ποσοστό των χρηστών που αναφέρουν ένα ύποπτο μήνυμα— σε συνδυασμό με τον \textbf{χρόνο έως την πρώτη αναφορά}. Μία μόνο έγκαιρη αναφορά επιτρέπει στην ομάδα ασφάλειας να αφαιρέσει ένα κακόβουλο μήνυμα από όλα τα γραμματοκιβώτια πριν το ανοίξουν οι περισσότεροι παραλήπτες.

Η κουλτούρα καθορίζει αν οι άνθρωποι θα προβούν σε αναφορά. Αν οι εργαζόμενοι που πατούν σε έναν σύνδεσμο ψαρέματος εκτίθενται δημόσια ή τιμωρούνται, θα αποκρύπτουν τα λάθη τους και ο οργανισμός θα χάνει την πιο έγκαιρη προειδοποίησή του. Μια κουλτούρα που ενθαρρύνει την αναφορά την καθιστά εύκολη (ένα κουμπί στο πρόγραμμα ηλεκτρονικού ταχυδρομείου), ευχαριστεί τους ανθρώπους για κάθε αναφορά —ακόμη και για τους ψευδείς συναγερμούς— και αντιμετωπίζει τα λάθη ως ευκαιρίες μάθησης. Με αυτόν τον τρόπο οι χρήστες γίνονται ενεργό επίπεδο ανίχνευσης και όχι αδυναμία του οργανισμού.

\section*{Βασικοί όροι}
\begin{description}[style=nextline,leftmargin=1.2em]
  \item[Κοινωνική μηχανική] Χειραγώγηση ανθρώπων ώστε να εκτελέσουν ενέργειες ή να αποκαλύψουν πληροφορίες προς όφελος του επιτιθέμενου.
  \item[Στοχευμένο ψάρεμα] Ηλεκτρονικό ψάρεμα προσαρμοσμένο σε συγκεκριμένο πρόσωπο ή ομάδα με στοιχεία που συλλέχθηκαν με έρευνα.
  \item[Κατασκευή προσχήματος] Επινόηση αληθοφανούς σεναρίου για την απόκτηση πληροφοριών ή πρόσβασης.
  \item[BEC] Παραβίαση επιχειρηματικού email: απάτη με υπόδυση αξιόπιστων προσώπων για εκτροπή πληρωμών.
  \item[Κόπωση MFA] Κατάκλυση ενός χρήστη με αιτήματα επαλήθευσης έως ότου εγκρίνει κάποιο.
  \item[DMARC] Πολιτική email που ορίζει στους παραλήπτες πώς να χειρίζονται μηνύματα που αποτυγχάνουν στους ελέγχους SPF/DKIM.
\end{description}

\section*{Σύνοψη κεφαλαίου}
\begin{itemize}
  \item Η κοινωνική μηχανική εκμεταλλεύεται αυτόματες αντιδράσεις απέναντι στην αυθεντία, το επείγον, την κοινωνική απόδειξη, τη συμπάθεια και την αμοιβαιότητα.
  \item Το ηλεκτρονικό ψάρεμα εμφανίζεται σε πολλές μορφές και κανάλια· η κατασκευή προσχήματος και το δόλωμα δημιουργούν καταστάσεις όπου οι επιβλαβείς ενέργειες μοιάζουν φυσιολογικές.
  \item Η φυσική πρόσβαση μπορεί να παρακάμψει τα δικτυακά μέτρα και πρέπει να προστατεύεται τόσο με τεχνολογία όσο και με τη συμπεριφορά του προσωπικού.
  \item Η BEC και η κόπωση MFA αντιμετωπίζονται με διαδικασίες επαλήθευσης, διπλή έγκριση, έλεγχο αυθεντικότητας email και MFA ανθεκτικό στο ψάρεμα.
  \item Ο κρίσιμος δείκτης είναι το ποσοστό αναφοράς και όχι το ποσοστό κλικ· μια κουλτούρα χωρίς απόδοση ευθυνών μετατρέπει τους χρήστες σε επίπεδο ανίχνευσης.
\end{itemize}

\section*{Ερωτήσεις ανασκόπησης}
\begin{enumerate}
  \item Εντοπίστε τις αρχές πειθούς που χρησιμοποιούνται στο μήνυμα: «Είμαι ο Διευθύνων Σύμβουλος. Χρειάζομαι την πληρωμή του προμηθευτή μέσα στα επόμενα 30 λεπτά —μην τηλεφωνήσετε, είμαι σε σύσκεψη».
  \item Γιατί η BEC δύσκολα σταματά με λογισμικό προστασίας από ιούς ή με περιβάλλοντα απομονωμένης εκτέλεσης; Ποια μέτρα είναι αποτελεσματικά;
  \item Εξηγήστε πώς η αντιστοίχιση αριθμού αποτρέπει μια βασική επίθεση κόπωσης MFA.
  \item Τεκμηριώστε γιατί το ποσοστό αναφοράς είναι καλύτερος δείκτης ευαισθητοποίησης από το ποσοστό κλικ.
\end{enumerate}
